\section{Set-theoretic preliminaries}\label{sec:preliminaries}

For undefined set-theoretic notions, we refer the reader to \cite{Kunen2009} and \cite{Kunen2011}.

The Axiom of Multiple Choice $\MC$ is the statement that
for every family of nonempty sets $(A_i)_{i\in I}$ there exists
a function $f$ whose domain is $I$ such that for every $i\in I$,
$f(i)$ is a nonempty finite subset of $A_i$.

For a set $Y$, let $h(Y)$ denote the Hartogs number
of $Y$, that is, the least
ordinal that does not inject into $Y$. Hartogs's theorem says that such a
number exists (see, e.g., \cite[\S I.11]{Kunen2009}).

We shall use that a finite union of finite sets is finite and that
choice functions for finite collections of nonempty sets exist in $\ZFm$.

The following is known as L\'evy's lemma \cite{Levy1962}.
See also
\cite[Lemma~1, p.~105]{Tachtsis2019}.

\begin{lemma}[L\'evy]\label{lem:partition}
In $\ZFm$, $\MC$ is equivalent to the statement that
every set has a partition into finite nonempty sets
indexed by an ordinal.

That is, for every set $X$ there exists an ordinal $\alpha$ and a family $(X_\beta)_{\beta<\alpha}$ of nonempty finite subsets of $X$ such that for every two distinct $\beta,\gamma<\alpha$, $X_\beta\cap X_\gamma=\varnothing$ and $\bigcup_{\beta<\alpha}X_\beta=X$.
\end{lemma}

